\section{Compositions and closures}

We have now constructed the eight possibilistic operators in formal concept analysis, as studied in \cite{dubois-dupin-prade_2007, dubois-prade_2009}, by standard categorical constructions associated to the underlying boolean profunctor $R\: G \nrightarrow G$. It remains then, to study the different possible compositions of upper-case and lower-case possibilistic operators. 

Given our eight operators, there are $16$ possible compositions that produce an endofunctor on $[G\op, \2]$. As we only consider endofunctors on $[G\op, \2]$ -- the ones on $[M,\2]\op$ are completely analogous -- we will use the combined notation $(-)_\sigma (-)^\rho := (-)^\sigma_\rho$ for simplicity. We then get the following table of compositions:

\begin{table}[ht]
\centering
\[
\begin{array}{c|cccc}
 & (-)_\Delta & (-)_\Pi & (-)_N & (-)_\nabla \\ \hline
(-)^\Delta
    & (-)^\Delta_\Delta
    & (-)^\Delta_\Pi
    & (-)^\Delta_N
    & (-)^\Delta_\nabla \\[0.4em]
(-)^\Pi
    & (-)^\Pi_\Delta
    & (-)^\Pi_\Pi
    & (-)^\Pi_N
    & (-)^\Pi_\nabla \\[0.4em]
(-)^N
    & (-)^N_\Delta
    & (-)^N_\Pi
    & (-)^N_N
    & (-)^N_\nabla \\[0.4em]
(-)^\nabla
    & (-)^\nabla_\Delta
    & (-)^\nabla_\Pi
    & (-)^\nabla_N
    & (-)^\nabla_\nabla
\end{array}
\]
\caption{The sixteen possible binary compositions.}
\label{tab:possibilistic-compositions}
\end{table}

The goal of this section is to prove that only three of the $16$ compositions that form closure operators on $P(G)$. First, we introduce some theory that allows us to prove this. 

\subsection{Idempotency and closure}

The standard FCA operator $(-)^\Delta_\Delta$ is a so-called closure operator, defined as follows. 

\begin{definition}
    A functor $F\: \C\to \C$ on a poset $\C$ is a \emph{closure operator} if it is extensive, monotone and idempotent, meaning it satisfies the following axioms:
\begin{enumerate}
    \item $A\leq F(A)$,
    \item $A\leq A' \implies F(A)\leq F(A')$, and
    \item $FF(A) = F(A)$.
\end{enumerate}
\end{definition}

The FCA adjunction coming from a formal context $(G,M,R)$ satisfies all three of these: by \cite[2.1.11]{jarvis_2025}, the monad $R_*R^*$ is idempotent, and the two other properties follow formally by the adjunction unit and the fact that it is a functor on posets. Thus, on a poset, any monadic functor determines a closure operator, and conversely, any closure operator is the monad of an adjunction -- see \cite[Chapter 7]{davey-priestley_2002} for details. 

In order to be recognized as a closure operator on \emph{some} formal context, these three axioms, then, have to be fulfilled. Otherwise the resulting fixedpoints have no chance of being formal concepts, or concepts of any type. In light of this discussion, the following proposition is rather trivial, but we state it for reference later.

\begin{proposition}
    \label{prop:axioms_for_closure}
    Let $L\: \C \to \C$ be an endofunctor on a poset. If $L$ fails to be extensive, monotone or idempotent, then it cannot be the closure operator of a formal context. 
\end{proposition}

This result will help us quickly identify which of the asymmetric compositions later possibly could give rise to alternative formal concepts on some related formal context. However, satisfying these three properties does not, as far as we know, guarantee that the fixedpoints form the formal concept of some formal context. 

\begin{remark}
    Dually, given an adjoint pair $L\dashv R$, the composition $LR$ is a comonad. On posets, this comonad satisfies dual properties of the closure operators, called being an \emph{interior} operator. In the case of formal concepts, instead of completing a subset $A$ to a formal concept, i.e., constructing the smallest formal concept containing the set, the interior operator finds the largest formal concept \emph{inside} $A$. 
\end{remark}


\subsection{Binary compositions}

We can now classify exactly which binary compositions of the eight operators that give closure operators. 

\begin{theorem}
    Out of the 16 operators in \cref{tab:possibilistic-compositions}, only $(-)^\Delta_\Delta$ and $(-)^\Pi_N$ are closure operators. 
\end{theorem}
\begin{proof}
    The monads of the adjunctions $(-)^\Delta \dashv (-)_\Delta$ and $(-)^\Pi \dashv (-)_N$, i.e., the operators $(-)^\Delta_\Delta$ and $(-)^\Pi_N$ are both closure operators by \cite[2.1.11]{jarvis_2025}. Similarly, due to the reversal of the adjunctions, $(-)^\Pi_N$ and $(-)^\nabla_\nabla$ are interior operators -- not closure operators. Hence, we only need to check the remaining $12$ cases. We check these by counterexamples, using \cref{prop:axioms_for_closure} to test whether $F$ satisfies the properties. In fact, it will be enough to test whether $F$ is extensive. 

    First, let $G=\{a,b\}$, $M=\{x,y\}$ and $R$ given by $aRx, bRy$. Chose the subset $A=G$. Then we have $A^\Delta = \varnothing$, $A^\Pi = \{x,y\}$, $A^N=\varnothing$ and $A^\nabla = \varnothing$. In this example, both $A^\Delta_\Pi = \varnothing$ and $A^\Delta_N = \varnothing$, showing that neither operators $(-)^\Delta_\Pi$ or $(-)^\Delta_N$ are closure operators. Similarly, we get $A^\Pi_\Delta = \varnothing$, $A^\Pi_\nabla=\varnothing$, $A^N_\Delta = \varnothing$, $A^N_\nabla = \varnothing$, $A^\Delta_\Pi = \varnothing$ and $A^\nabla_N = \varnothing$, showing that none of these can be closure operators as well. 

    It remains to show the same for $(-)^\Pi_\Pi$, $(-)^\Delta_\nabla$, $(-)^N_N$ and $(-)^\nabla_\Delta$. For the former two we use the formal context $G=\{a\}$, $M=\varnothing$ and $A=G$. This gives $A^\Pi_\Pi = \varnothing$ and $A^\Delta_\nabla = \varnothing$, showing neither is extensive and hence not a closure operator. Similarly, using $G=\{a\}$, $M=\{x\}$, $aRx$ and $A=G$ $A^\nabla_\Delta = \varnothing$. For the final operator, $(-)^N_N$ we use the context $G=\{a,b\}$, $M=\{x\}$, $aRx$, $bRx$ and $A=\{a\}$. This gives $A^N_N = \varnothing$, which also finishes the proof. 
\end{proof}

\subsection{Other closure operations}

We round off this section by using the categorical machinery to identify a sequence of new types of concepts, or pairs, that naturally arise from the framework. These come from the following observation: Given two adjunctions 
\begin{center}
\begin{tikzcd}
    \C \arrow[r, yshift=2pt, "F_1"] & \D \arrow[r, yshift=2pt, "F_2"] \arrow[l, yshift=-2pt, "G_1"] & \E \arrow[l,yshift=-2pt, "G_2"]
\end{tikzcd}
\end{center}
the composed functor $F_2 F_1$ is a left adjoint to $G_2 G_1$. Recall from \cref{prop:axioms_for_closure} that on complete lattices, closure operators are equivalent to monads. This gives the following result. 

\begin{theorem}
    The functors $(-)^\Delta$ and $(-)^\Pi$ on $P(G)$, and  similarly $(-)_\Pi$ and $(-)_\nabla$ on $P(M)\op$ are left adjoints. Any composition of these will produce a closure operator on $P(G)$ by composing with the corresponding right adjoint. 
\end{theorem}

\begin{example}
    \label{ex:new-operator}
    The functor 
    \[(-)^\Delta (-)_\Pi (-)^\Delta \: [G\op, \2]\to [M,\2]\op\]
    is a composition of left adjoints, hence is a left adjoint itself, with right adjoint given by $(-)_\Delta (-)^N (-)_\Delta$. The composition of these is by \cite[2.1.11]{jarvis_2025} and standard categorical arguments a closure operator. 
\end{example}

\begin{proposition}
    If $C$ is the closure operator defined in \cref{ex:new-operator}, then there is no formal context $(G,M,S)$ such that $C = (-)_{\Delta_S}(-)^{\Delta_S}$. 
\end{proposition}
\begin{proof}
    Any FCA closure operator can be computed poitwise via $A^\Delta_\Delta = \bigcap_{g\in A}g^\Delta_\Delta$. The operator $C$ does not always satisfy this property, as we can see by the counterexample $G=\{a,b\}$, $M=\{x,y\}$ with $aRy$, $bRx$. Here $C(G)=G$, while $C(\{a\})\cap C(\{b\})=\varnothing$. 
\end{proof}

Another nice feature of the categorical language is that one can prove the existence of other adjoints, which are not one of the eight possibilistic operators. For example, as $(-)_\Delta$ is a right Kan extension, it preserves limits. As $[M,\2]\op$ is a complete lattice, the adjoint functor theorem applies, meaning it has a \emph{left} adjoint 
\[L^\Delta\: [G\op,\2]\to [M,\2]\op.\] 
This is not one of the possibilistic operators, and does not coincide with mixing quantifiers with the relation $R$ and its complement $\oR$. The composition $(-)_\Delta L^\Delta$ is a closure operator on $P(G)$ due to the general property of being an adjoint pair on posets. 

Similarly, the functor $(-)_\nabla$ is the dual of a left Kan extension, hence preserve limits, meaning that it has a left adjoint $L^\nabla$ by the adjoint functor theorem. The composition $(-)_\nabla L^\nabla$ is, by the adjunction property, another closure operator on $P(G)$. 